\section{Особенности реализации}

Программа написана на языке Go. Функции преобразования находятся в модуле
\texttt{bf16}. В них реализовано кодирование десятичного
числа в 16-битный формат \texttt{bfloat16}, обратное восстановление числа из
строки битов и вывод 16 бит в строку. Формат состоит из одного бита
знака, восьми битов смещённой экспоненты и семи битов мантиссы.

\vspace{8pt}
\subsection*{Функция FloatToBits}
\textbf{Вход:} \texttt{float64 value} --- десятичное вещественное число. \par
\textbf{Выход:} \texttt{uint16} --- 16-битное представление числа в формате
\texttt{bfloat16}. \par
\textbf{Описание:} Функция определяет знак числа, отдельно переводит его
целую и дробную части в двоичный вид, а затем нормализует полученную запись.
После этого формируются семь бит мантиссы, экспонента смещается на 127, и
поля $S$, $E$, $M$ собираются сдвигами. При переполнении возвращается
бесконечность, при слишком маленькой экспоненте --- ноль. \par
\vspace{6pt}

\begin{lstlisting}[style=common, language=Go, caption=FloatToBits, label=lst:float-to-bits]
func FloatToBits(value float64) uint16 {
    var sign uint16
    if value < 0 {
        sign, value = 1, -value
    }
    if value != value {
        return nanBits
    }
    if value == 0 {
        return sign << 15
    }

    integer := uint64(value)
    integerBits := integerToBinary(integer)
    fractionBits := fractionToBinary(value - float64(integer))
    exponent, source, found := normalize(integerBits, fractionBits)
    if !found {
        return sign << 15
    }
    mantissa, bump := roundMantissa(source)
    exponent += bump

    exponent += bias
    if exponent >= maxExponent {
        return sign<<15 | infinityBits
    }
    if exponent <= 0 {
        return sign << 15
    }
    return sign<<15 | uint16(exponent<<7) | mantissa
}
\end{lstlisting}

\newpage

\subsection*{Функция BitsToBinary}
\textbf{Вход:} \texttt{uint16 bits} --- 16-битное число. \par
\textbf{Выход:} \texttt{string} --- строка из 16 символов \texttt{0} и
\texttt{1}. \par
\textbf{Описание:} Функция вызывает общий перевод \texttt{toBinary} с флагом
\texttt{true}. Поэтому результат всегда содержит 16 бит и нули слева.

\begin{lstlisting}[style=common, language=Go, caption=BitsToBinary, label=lst:bits-to-binary]
func BitsToBinary(bits uint16) string {
    return toBinary(uint64(bits), true)
}
\end{lstlisting}


\subsection*{Функция integerToBinary} 
\textbf{Вход:} \texttt{uint64 value} --- неотрицательная целая часть числа. \par
\textbf{Выход:} \texttt{string} --- двоичная запись числа. \par
\textbf{Описание:} Функция вызывает общий перевод \texttt{toBinary} без
добавления нулей слева. Поэтому результат содержит только значимые биты.

\begin{lstlisting}[style=common, language=Go, caption=integerToBinary, label=lst:integer-to-binary]
func integerToBinary(value uint64) string {
    return toBinary(value, false)
}

\end{lstlisting}

\subsection*{Функция toBinary}
\textbf{Вход:} \texttt{uint64 value} --- число для перевода;
\texttt{bool fillTo16} --- флаг добавления нулей до 16 бит. \par
\textbf{Выход:} \texttt{string} --- двоичная запись числа. \par
\textbf{Описание:} Функция получает младший бит, добавляет его в начало
строки и сдвигает число вправо. Если \texttt{fillTo16} равен \texttt{true},
слева добавляются нули до длины 16 символов.

\begin{lstlisting}[style=common, language=Go, caption=toBinary, label=lst:to-binary]
func toBinary(value uint64, fillTo16 bool) string {
    bits := ""
    for {
        if value&1 == 1 {
            bits = "1" + bits
        } else {
            bits = "0" + bits
        }
        value >>= 1
        if value == 0 {
            break
        }
    }
    if fillTo16 {
        for len(bits) < totalBits {
            bits = "0" + bits
        }
    }
    return bits
}
\end{lstlisting}

\subsection*{Функция fractionToBinary}
\textbf{Вход:} \texttt{float64 value} --- дробная часть числа от 0 до 1. \par
\textbf{Выход:} \texttt{string} --- двоичные разряды дробной части. \par
\textbf{Описание:} Дробная часть последовательно умножается на 2. Целая
часть результата становится очередным битом, после чего она вычитается.
Строится не более 135 бит: этого достаточно для определения мантиссы и
экспоненты \texttt{bfloat16}.

\begin{lstlisting}[style=common, language=Go, caption=fractionToBinary, label=lst:fraction-to-binary]
func fractionToBinary(value float64) string {
    bits := ""
    for index := 0; value != 0 && index < 135; index++ {
        value *= 2
        if value >= 1 {
            bits += "1"
            value--
        } else {
            bits += "0"
        }
    }
    return bits
}
\end{lstlisting}

\subsection*{Функция normalize}
\textbf{Вход:} \texttt{string integerBits}, \texttt{string fractionBits} ---
двоичные записи целой и дробной частей. \par
\textbf{Выход:} истинная экспонента, исходные биты мантиссы и флаг успешной
нормализации. \par
\textbf{Описание:} Если целая часть не равна нулю, её старшая единица становится
неявной единицей перед точкой, а экспонента равна позиции этой единицы.
Если целая часть нулевая, функция ищет первую единицу в дробной части;
экспонента в этом случае отрицательная.

\begin{lstlisting}[style=common, language=Go, caption=normalize, label=lst:normalize]
func normalize(integerBits, fractionBits string) (int, string, bool) {
    if integerBits != "0" {
        return len(integerBits) - 1, integerBits[1:] + fractionBits, true
    }

    for index, bit := range fractionBits {
        if bit == '1' {
            return -(index + 1), fractionBits[index+1:], true
        }
    }
    return 0, "", false
}
\end{lstlisting}

\newpage

\subsection*{Функция roundMantissa}
\textbf{Вход:} \texttt{string source} --- биты после неявной единицы. \par
\textbf{Выход:} 7-битная мантисса и флаг переноса в экспоненту. \par
\textbf{Описание:} Сохраняются первые семь бит. Восьмой бит определяет
округление: если он равен единице, к мантиссе прибавляется единица. При
переполнении мантисса обнуляется, а флаг переноса равен единице.

\begin{lstlisting}[style=common, language=Go, caption=roundMantissa, label=lst:round-mantissa]
func roundMantissa(source string) (uint16, int) {
    for len(source) < 8 {
        source += "0"
    }

    var mantissa uint16
    for _, bit := range source[:7] {
        mantissa <<= 1
        if bit == '1' {
            mantissa |= 1
        }
    }
    if source[7] == '1' {
        mantissa++
    }
    if mantissa == 128 {
        return 0, 1
    }
    return mantissa, 0
}
\end{lstlisting}

\subsection*{Функция BitsToFloat}
\textbf{Вход:} \texttt{string binary} --- строка из 16 бит
\texttt{bfloat16}. \par
\textbf{Выход:} \texttt{float64} --- восстановленное десятичное число. \par
\textbf{Описание:} Из строки отдельно берутся знак, восемь бит экспоненты
и семь бит мантиссы. Для обычного числа
используется формула $(-1)^S(1+M)2^{E-127}$, степень двойки вычисляется
функцией \texttt{math.Pow}. При нулевой экспоненте
восстанавливается денормализованное число, а экспонента из единиц обозначает
бесконечность или \texttt{NaN}.

\begin{lstlisting}[style=common, language=Go, caption=BitsToFloat, label=lst:bits-to-float]
func BitsToFloat(binary string) float64 {
    sign := 1.0
    if binary[0] == '1' {
        sign = -1
    }
    exponent := binaryToNumber(binary[1 : 1+exponentBits])
    mantissa := bitsToFraction(binary[1+exponentBits:])

    if exponent == maxExponent {
        zero := 0.0
        if mantissa != 0 {
            return zero / zero
        }
        return sign / zero
    }

    value, power := mantissa, 1-bias
    if exponent != 0 {
        value, power = 1+mantissa, int(exponent)-bias
    }
    return sign * value * math.Pow(2, float64(power))
}
\end{lstlisting}

\subsection*{Функция binaryToNumber}
\textbf{Вход:} \texttt{string binary} --- строка из нулей и единиц. \par
\textbf{Выход:} \texttt{int} --- число в десятичной системе. \par
\textbf{Описание:} Функция проходит по битам слева направо. На каждом шаге
накопленное число сдвигается влево, а при единице увеличивается на единицу.

\begin{lstlisting}[style=common, language=Go, caption=binaryToNumber, label=lst:binary-to-number]
func binaryToNumber(binary string) int {
    value := 0
    for _, bit := range binary {
        value <<= 1
        if bit == '1' {
            value++
        }
    }
    return value
}
\end{lstlisting}

\subsection*{Функция bitsToFraction}
\textbf{Вход:} \texttt{string binary} --- биты мантиссы. \par
\textbf{Выход:} \texttt{float64} --- дробная часть мантиссы. \par
\textbf{Описание:} Функция рассматривает первый бит как $2^{-1}$, второй как
$2^{-2}$ и так далее. Степень вычисляется функцией \texttt{math.Pow}; если
бит равен единице, его вес добавляется к результату.

\begin{lstlisting}[style=common, language=Go, caption=bitsToFraction, label=lst:bits-to-fraction]
func bitsToFraction(binary string) float64 {
    fraction := 0.0
    for index, bit := range binary {
        if bit == '1' {
            fraction += math.Pow(2, -float64(index+1))
        }
    }
    return fraction
}
\end{lstlisting}

\newpage
